\Needspace{6\baselineskip}\SourceLeaf{86}

\[
= X^3\frac{d^2\varphi x}{dt\,du} + 3X^2\frac{dX}{dt}\frac{d\varphi x}{dt}
= \frac{dX^3 \frac{d\varphi x}{dt}}{dt},
\]

\noindent
it follows that

\begin{equation}\tag{5}
\frac{d^3\varphi x}{du^3} = \frac{d^2X^3 \frac{d\varphi x}{dt}}{dt^2}.
\end{equation}

\noindent\hspace*{\parindent}One likewise finds

\[
\frac{d^4\varphi x}{du^4} = \frac{d^3X^4 \frac{d\varphi x}{dt}}{dt^3} \quad \text{u.\,s.\,f.}
\]

\Needspace{9\baselineskip}\noindent\hspace*{\parindent}\textls{In general, however}, differentiating $X^n\dfrac{d\varphi x}{dt}$ with respect to $u$ and applying
substitutions (3) and (4) gives

\begin{equation}\tag{6}
\frac{dX^n \frac{d\varphi x}{dt}}{du} = \frac{dX^{n+1} \frac{d\varphi x}{dt}}{dt}.
\end{equation}

\noindent
If, then, the law holds through term $n^{\text{te}}$, or

\[
\frac{d^n\varphi x}{du^n} = \frac{d^{n-1}X^n \frac{d\varphi x}{dt}}{dt^{n-1}},
\]

\noindent
then differentiating with respect to $u$ gives

\[
\frac{d^{n+1}\varphi x}{du^{n+1}} = \frac{d\,d^{n-1}X^n\frac{d\varphi x}{dt}}{du . dt^{n-1}}
= \frac{d^{n-1}dX^n\frac{d\varphi x}{dt}}{dt^{n-1}du},
\]

\noindent
[thus], by (6),

\begin{equation}\tag{7}
\frac{d^{n+1}\varphi x}{du^{n+1}} = \frac{d^nX^{n+1} \frac{d\varphi x}{dt}}{dt^n},
\end{equation}

\noindent
i.\,e.\ [the rule holds] for the next term as well.

\noindent\hspace*{\parindent}Thus equations (1), (4), (5) give the transformations above for the first three
terms, and (6) establishes their validity for all the following terms.

\medskip

\Needspace{10\baselineskip}\begin{center}
\textbf{REMARK ON THE PAPER »NEW PROOF OF LAGRANGE'S THEOREM«.}
\end{center}

\noindent\hspace*{\parindent}According to an entry in the diary mentioned on p. 20, \textsc{Gauss} found a
proof of \textsc{Lagrange}'s theorem on 27 December 1796. The paper printed above is
written on a single slip of paper without a date, but there can be no doubt that this is
the derivation \textsc{Gauss} found at the end of 1796. Publication at the time was
prevented only by various chance circumstances
(cf.\ \textsc{Sartorius von Waltershausen}, Gauss zum Gedächtniss, p. 22).

\begin{flushright}
\textsc{Fricke}.
\end{flushright}
